\section{The literal base-change-to-Speh bridge}
\label{sec:base-change-v44}

This chapter proves the complete finite consumer
\[
 \mathrm{BaseChangeIsobaricN}\ \Longrightarrow\
 \mathrm{DefiniteN}\ \Longrightarrow\ \mathrm{BcSpeh}.
 \tag{BCS.1}
\]
The first premise is the original named hypothesis
\texttt{H.litN.aut.bc}.  It is retained with its full definite-inner-form,
all-rank and nongeneric scope.  The chapter does not supply a new
trace-formula proof of that premise.  The other named inputs retained
from the original literature record are
\texttt{H.litN.loc.bz}, \texttt{bern}, \texttt{sq}, \texttt{cos},
\texttt{H.litN.aut.fin}, and the boundedness clause of
\texttt{H.litN.aut.cusp}.  A precisely stated published local
generic-unitary decomposition is matched in BCS8.  Basic Haar measure,
the finite-volume general-linear adelic quotient, Fourier uniqueness on
compact additive adelic quotients, and highest-weight theory for
\(U(n)\) are the original analytic and representation-theoretic
foundations.  Their uses are identified below.  None of these inputs
is silently counted as a newly reconstructed theorem.

Let \(E/E^+\) be the original CM number field.  It is unramified at
all finite places, as required by the base-change hypothesis.
Let \(n\ge1\), and
let \(\pi\) be an actual irreducible subspace of the original
automorphic functions on the definite unitary adelic group
\(U_{\mathbb A}\), with \(\mathrm{IsAutRepU}(\pi)\).
At a finite place \(w\), put \(F=E_w\),
\(K_r=\mathrm{GL}_r(\mathcal O_F)\), and
\(q_w=\#(\mathcal O_E/w)\).  All products \(\times\) below are
normalized induction from the upper standard block parabolic, with
the original increasing block order.  Thus
\[
 f(pg)=\delta_P(p)^{1/2}\rho(m_p)f(g),\qquad
 \nu^t(g)=|\det g|_F^t,\qquad
 s_{d,k}=\frac{d-1}{2}-k\quad(0\le k<d).
 \tag{BCS.2}
\]
For a multiset \(m\), write
\[
 \operatorname{spehInf}(d,m)
   =\sum_{k=0}^{d-1}\{z+s_{d,k}:z\in m\}.
 \tag{BCS.3}
\]
The definition of \(\mathrm{IsSpeh}(\rho,d,\sigma)\) is stronger
than being an unspecified irreducible quotient: \(\sigma\) is the
unique irreducible cosocle of
\(\rho\nu^{s_{d,0}}\times\cdots\times\rho\nu^{s_{d,d-1}}\),
and every maximal proper invariant subspace of that induction is the
kernel of its quotient map.  We prove this precise assertion.

\subsection{Replacing the selected subquotient by the full local product}

% BEGIN BCS1
\paragraph{BCS1. The actual discrete and normalized-cuspidal local forms.}
\label{claim:BCS1}
Every local component of an original discrete automorphic
\(\mathrm{GL}_r\)-space is unitarisable.  A normalized cuspidal
space is itself discrete, so the same conclusion holds for each of
its local components.
\begin{proof}
The original \(\mathrm{IsDiscrete}(V)\) includes square
integrability of every actual function in \(V\), invariance under
\(\mathrm{GL}_r(E)\) and the real split center, stability under all
finite, compact-infinity and Lie operators, and nonzero irreducibility.
Choose one Haar measure and one fundamental domain for the source
subgroup
\(\mathrm{GL}_r(E)\langle\operatorname{splitCentre}(1)\rangle\).
The definition permits separate choices for separate functions; Haar
uniqueness and invariance of the integral under changing a fundamental
domain put every function in \(L^2\) for this common choice.  Polarization
and Cauchy--Schwarz therefore define a single form
\[
 B(f,h)=\int_{\mathcal F}\overline{f(g)}h(g)\,dg.
 \tag{BCS.4}
\]
It is positive definite on the actual continuous functions.  Indeed a
nonzero value gives an open set of positive quotient Haar measure on
which the absolute value is bounded below.  Such a function cannot have
zero \(L^2\) norm.  The adelic general linear group is unimodular;
right translation preserves the quotient measure and hence (BCS.4).

An original local component is a nonzero local intertwining map
\(T:\sigma\to V\).  Its kernel is invariant and \(\sigma\) is
irreducible, so \(T\) is injective.  The pullback \(B(Tv,Tz)\) is
the required positive invariant Hermitian form on \(\sigma\).
No completion of its underlying algebraic vector space is necessary.

For a normalized cusp \(P\), the normalization says that every
function is fixed by the same real split center.  The retained
\texttt{CuspFormsBoundedGL} clause bounds each such cusp function.
The original quotient just used has finite volume, so boundedness
implies its square integrability.  The other clauses of
\(\mathrm{IsDiscrete}(P.V)\) are already clauses of the cuspidal
subspace.  Apply the preceding pullback argument to \(P.V\).
This last argument uses only normalization, not an unproved assertion
that a cuspidal local component is tempered.
\end{proof}
% END BCS1

% BEGIN BCS2
\paragraph{BCS2. All finite unitary local products, with the original order.}
\label{claim:BCS2}
Let \(\sigma_i\) be smooth irreducible unitarisable representations
of \(\mathrm{GL}_{N_i}(F)\), with \(N_i>0\), and let
\(\sum_iN_i=n\).  Their original normalized product
\(\operatorname{locProd}(h,w,\sigma)\) is smooth, unitarisable and
irreducible.
\begin{proof}
Tensor the positive forms on the inducing spaces.  By Iwasawa
decomposition, normalized induction is its space of locally constant
sections over the compact quotient \((P\cap K_n)\backslash K_n\).
The formula
\[
             (f_1,f_2)=\int_{K_n}B(f_1(k),f_2(k))\,dk,
             \qquad\operatorname{vol}(K_n)=1,
 \tag{BCS.5}
\]
is positive definite.  It is invariant under right translation by
\(G\): the change of Iwasawa coordinates has Jacobian
\(\delta_P^{-1}\), while the two covariance factors contribute
\(\delta_P\).  Smoothness follows from the right open stabilizer
in the induced model.  Thus normalized induction preserves
unitarisability in precisely this normalization.

For two factors, irreducibility is the retained literal
\texttt{BernsteinUnitaryProductIrreducible} hypothesis.  Induct on
the number of factors.  The product of the first \(c-1\) factors
is smooth irreducible and unitarisable by the induction hypothesis,
so the same two-factor theorem applies to it and the last factor.
Induction in stages identifies this result with the original full
product.  Explicitly, a section is sent to the function of two group
variables obtained by successive evaluation; the inverse evaluates
the inner section.  The modulus identity
\(\delta_P^{1/2}=\delta_Q^{1/2}\delta_{P\cap M_Q}^{1/2}\)
checks covariance.  Both maps use the original block embeddings and
their original order.  One factor gives the representation itself;
the empty product, when needed at rank zero, is the one-dimensional
representation of the trivial group.  No distinctness of the factors
and no bound on their number enters this proof.
\end{proof}
% END BCS2

% BEGIN BCS3
\paragraph{BCS3. The full original \texttt{DefiniteN} bridge.}
\label{claim:BCS3}
The original \texttt{BaseChangeIsobaricN} and the inputs of BCS1--BCS2
imply the entire original \texttt{DefiniteN}.
\begin{proof}
Take exactly the base-change witnesses
\(c,N_i,h,V_i,\sigma_{i,w},M_{i,\tau},Q_w\).
Here \(Q_w\) is the selected irreducible subquotient of the full
product \(I_w=\operatorname{locProd}(h,w,\sigma_{\bullet,w})\).
BCS1 makes every \(\sigma_{i,w}\) unitarisable, and BCS2 makes
\(I_w\) irreducible.  A nonzero irreducible subquotient of an
irreducible representation is isomorphic to it: the invariant
subspace from which the quotient is taken is nonzero, hence the
whole space; the quotient kernel is proper, hence zero.
Choose the resulting equivalence \(e_w:Q_w\simeq I_w\).

Every fixed vector and every finite Hecke sum transports along
\(e_w\).  Thus the supplied spherical clause at every
\(\mathrm{IsUnramU}(\pi,w)\) becomes the spherical clause for
\(I_w\).  For every permitted \(w,s,\varpi,q,\alpha\), the
supplied inert clause becomes
\[
 \mathrm{HasSatGL}\bigl(w,I_w,\operatorname{bcParam}(n,\alpha)\bigr).
 \tag{BCS.6}
\]
This keeps the original conditions that \(s\) extends CM conjugation,
that \(\mathrm{IsUnramInv}(\varpi,s,q)\) holds, and that
\(\mathrm{HasSatAt}(\pi,w,s,\varpi,q,\alpha)\) holds.

At a split place the original nonzero map from \(Q_w\) into the
actual unitary automorphic space is precomposed with \(e_w^{-1}\).
Its image is still in \(\pi\), and its intertwining identity is
still along
\((cLoc_w)\circ(\operatorname{locU}(n,w)\hookrightarrow U_{\mathbb A})\).
Hence it gives the full \(\mathrm{IsCompAlong}\) clause for
\(I_w\), with precisely that subgroup and map.  All other fields
---positive \(N_i\), discreteness, \(\mathrm{CsdAsc}\), actual
local components, infinitesimal characters, and the weight sum---are
copied without change.  These are exactly all fields of
\(\mathrm{IsBaseChangeU}\), proving \texttt{DefiniteN}.
\end{proof}
% END BCS3

\subsection{The bounded residual-spectrum consumer}

The input \(\mathrm{CsdAsc}(V)\) already supplies one normalized
cuspidal \(P\) of rank \(a\) and a length \(d\), with
\(da=N\), together with an injective map from \(V\) to the
original unnormalized induced automorphic functions.  The map
intertwines all finite, compact-infinity and Lie operators.  Its
\(k\)-th cuspidal coefficient is
\[
                P\,|\det|_{\mathbb A_E}^{(a-1)s_{d,k}}.
 \tag{BCS.7}
\]
The following argument consumes this actual embedding.  It does not
infer the embedding from a classification assertion.

% BEGIN BCS4
\paragraph{BCS4. Actual local components, isotypy and infinity characters.}
\label{claim:BCS4}
Every original irreducible cuspidal or discrete automorphic space
has an actual irreducible smooth local component at every finite
place and an infinitesimal-character multiset at every embedding.
Its restriction to a fixed finite local group is an algebraic sum
of copies of any one of its local components.  In particular all
such components are isomorphic.  The local existence assertion uses
the retained \texttt{AutomorphicFormsFiniteGL}.
\begin{proof}
First, every central Capelli operator preserves the finite-dimensional
space generated from a nonzero automorphic vector by the center.
It has an eigenvector.  Since it commutes with all three kinds of
automorphic operators, its eigenspace is a nonzero globally stable
subspace, hence the entire irreducible space.  The Capelli expression
\(C_r(u)=\operatorname{qdet}(E_{ij},u+(r-1)/2)\) is monic of
degree \(r\) in \(u\).  Evaluating on a fixed nonzero vector and
a linear functional taking value one gives a monic scalar polynomial.
Factoring it over \(\mathbb C\) gives a multiset \(m\) of size
\(r\) such that
\(C_r(u)=\prod_{z\in m}(u+z)\) on the entire space.  This is the
literal \(\mathrm{HasInfCharGL}\), or \(P.\mathrm{HasInfChar}\).
The commutation identity used here follows by expanding the Capelli
column determinant and using
\([E_{ij},E_{kl}]=\delta_{jk}E_{il}-\delta_{li}E_{kj}\);
the two terms associated with each off-diagonal replacement cancel,
including the consecutive diagonal shifts.  Equal scalar polynomials
have equal multisets, so these characters are unique on a nonzero
space and are preserved by all-operator equivalences.

For completeness, local existence can be obtained without assuming
a restricted tensor product in advance.  Fix a nonzero vector \(v\)
and a compact open local subgroup \(K\) fixing it.  Let \(e_K\)
be normalized compact averaging.  The local subspace
\[
                  Z=\bigcap_{g\in\mathrm{GL}_r(F)}
                       \ker(e_K\,g)
 \tag{BCS.8}
\]
is stable under the local group and every operator commuting with
it.  Finite adelic operators factor as a local operator and one
trivial at this place; infinity operators commute with it.  Thus
\(Z\) is globally stable.  Since \(e_Kv=v\ne0\), global
irreducibility gives \(Z=0\).  Consequently every nonzero local
subrepresentation has nonzero \(K\)-fixed vectors: otherwise
averaging every translate puts it inside \(Z\).

Take the nonzero local-stable space \(T\) of vectors with the
fixed compact-infinity type generated by \(v\), fixed by a chosen
open compact group away from the place.  It has the central
character just proved.  This finite-dimensional infinity type is
continuous: evaluations separate its finite span, so finitely many
evaluations give coordinates, and each coordinate of its orbit is
continuous because the original automorphic functions are smooth.
For each local compact open \(K'\), its
\(K'\)-fixed part is finite dimensional by the retained automorphic
finiteness theorem.  This is an application to an actual open finite
adelic level: intersect the original level with the inverse image of
\(K'\) under the local projection and with the local inclusion
condition.  Thus all original level and infinity-type conditions are
preserved.

Among the nonzero invariant subspaces \(W\le T\), choose one
minimizing the positive integer \(\dim W^K\).  Let \(X\) be the
local span of \(W^K\).  A nonzero invariant subspace \(Y\le X\)
has \(\dim Y^K\ge\dim W^K\); since \(Y^K\subseteq W^K\),
they are equal.  Hence \(Y\) contains \(W^K\) and so equals
\(X\).  This gives an actual irreducible smooth subspace and its
inclusion, not just a composition factor.  It is admissible, either
by the same finite-dimensional fixed spaces or by
\texttt{BernsteinZelevinskyFiniteness}.

Now fix any irreducible local copy \(X\).  The sum of the images
of all local maps \(X\to V\) is stable under local operators and
under their commuting operators, hence is all of \(V\).  A sum of
simple submodules is semisimple: enlarge a direct sum of them until
it is maximal; any further simple submodule either meets it trivially
or is contained in it.  This proves algebraic isotypy.  Schur's
endomorphism field is \(\mathbb C\): an endomorphism restricts to
a nonzero finite-dimensional compact-fixed space, has an eigenvalue
there, and irreducibility makes its eigenspace the whole local
representation.  Thus a nonzero vector in any tensor of these
isotypic restrictions is detected by a local-equivariant projection
onto the tensor of the chosen components.  Such a tensor is simple,
as successive use of the density theorem for a simple module with
scalar endomorphisms separates the finitely many tensor coordinates.
\end{proof}
% END BCS4

% BEGIN BCS5
\paragraph{BCS5. Cuspidal local components are generic.}
\label{claim:BCS5}
Every actual finite local component of a positive-rank original cusp
admits a nonzero nondegenerate Whittaker functional.
\begin{proof}
Choose the standard continuous additive character of
\(\mathbb A_E/E\), nontrivial at each completion.  We recall why
cuspidality gives the needed nonzero coefficient.  Expand a nonzero
cusp function along the last-column additive radical.  The constant
coefficient is its proper-parabolic constant term and is zero.
Fourier uniqueness on the compact quotient
\((\mathbb A_E/E)^{a-1}\) leaves a nonzero nonconstant coefficient.
The rational group \(\mathrm{GL}_{a-1}(E)\) is transitive on
nonzero row characters, so rational left invariance identifies these
coefficients with translates of the one using the last coordinate.
That coordinate coefficient therefore cannot vanish identically.

Repeat on the upper-left \((a-1)\)-block.  At each stage the
previously imposed last-coordinate characters commute with the next
column integration.  The zero new coefficient is the constant term
along the corresponding larger proper parabolic, followed by the
previous character integrations; it vanishes by cuspidality and
Fubini on compact additive quotients.  The same rational transitivity
and Fourier uniqueness show that the next designated nonconstant
coefficient is nonzero as a function.  Induction ends with a
nonzero function
\[
 W_f(g)=\int_{N(E)\backslash N(\mathbb A_E)}
     f(ug)\,\psi\!\left(-\sum_{j=1}^{a-1}u_{j,j+1}\right)du.
 \tag{BCS.9}
\]
Successive integration in the column coordinates is a definition of
this integral; compactness makes all integrations finite.  This
argument also proves the transformation law under \(N(\mathbb A_E)\)
by translating the integration coordinates.  For rank one, the
unipotent group is trivial and \(W_f=f\).

Let \(T:\rho\to P.V\) be the given nonzero local map.  It is
injective, so choose \(v\) with \(Tv\ne0\).  Some
\(W_{Tv}(g_0)\) is nonzero.  Factor \(g_0=h g_w\), where \(h\)
has identity component at \(w\).  Then
\(\ell(x)=W_{Tx}(h)\) is nonzero on \(\rho(g_w)v\).
The element \(h\) commutes with the local group.  The transformation
law from (BCS.9) therefore gives
\[
 \ell(\rho(u)x)=\psi_w\!\left(\sum_j u_{j,j+1}\right)\ell(x),
                         \qquad u\in N(E_w).
\]
The local character is nontrivial.  This is exactly genericity of
the given local representation.
\end{proof}
% END BCS5

% BEGIN BCS6
\paragraph{BCS6. Smooth Hermitian duality turns the ascending map into a quotient.}
\label{claim:BCS6}
Let \(\rho,\sigma\) be smooth irreducible unitarisable
representations of \(\mathrm{GL}_a(F),\mathrm{GL}_{ad}(F)\), with
\(d\ge1\).  A nonzero map
\[
 \sigma\longrightarrow
       \rho\nu^{-s_{d,0}}\times\cdots\times\rho\nu^{-s_{d,d-1}}
 \tag{BCS.10}
\]
implies a surjective map from the descending product
\(\rho\nu^{s_{d,0}}\times\cdots\times\rho\nu^{s_{d,d-1}}\)
onto \(\sigma\).
\begin{proof}
For a compact open \(K\), normalized averaging is a finite average
on each smooth vector: choose an open normal subgroup of \(K\)
fixing that vector and average over the finite quotient.  It is an
idempotent onto fixed vectors.  In particular, the smooth dual of an
injection is surjective.  To extend a smooth functional on the smaller
space, first extend it algebraically, then average the extension over
a compact open subgroup fixing the original functional.  The
restriction remains the original functional.  The same argument
applies to the conjugate, or Hermitian, smooth dual \(V^h\).

For an admissible unitary representation with positive form \(B\),
the map \(v\mapsto B(v,\cdot)\), interpreted with the conjugate
scalar structure, identifies it with its Hermitian smooth dual.  To
see surjectivity, a \(K\)-fixed functional is determined by its
restriction to the finite-dimensional space \(V^K\); there the
positive form represents it.  Orthogonal averaging shows that the
same representing vector works on all of \(V\).  The needed
admissibility is the literal irreducible-admissibility clause of
\texttt{H.litN.loc.bz}.

The normalized-induction duality in that same record, paired by the
integral over \(K_{ad}\), identifies the Hermitian dual of an induced
product with the product of the Hermitian duals in the \emph{same}
upper-parabolic order.  Complex conjugation followed by ordinary
contragredience sends \(\rho\nu^t\) to
\(\rho^h\nu^{-t}\) for real \(t\).
The map (BCS.10) is injective by irreducibility.  Take its Hermitian
smooth dual, use the preceding surjectivity, and then the unitary
identifications \(\rho^h\simeq\rho\) and
\(\sigma^h\simeq\sigma\).  This gives the asserted quotient,
with each exponent \(-s_{d,k}\) changed to \(s_{d,k}\).
There is no reordering theorem or local Gelfand--Kazhdan assumption
hidden in this step.
\end{proof}
% END BCS6

% BEGIN BCS7
\paragraph{BCS7. Extracting the exact ascending local induction.}
\label{claim:BCS7}
\leavevmode\par
Suppose \(V\) is discrete, \(\operatorname{AscEmb}(V,P,d)\),
\(N\ge1\), and \(\rho,\sigma\) are any actual local
components of \(P,V\) at \(w\).  Then \(da=N\), both
\(a,d\) are positive, and there is a nonzero map (BCS.10).
\begin{proof}
The rank identity is the block identity in \(\operatorname{AscEmb}\).
Its positive product gives \(a,d\ge1\).  Compose the actual
component map \(\sigma\to V\) with the injective ascending
embedding to get a nonzero local map \(F\) into the given space
of induced automorphic functions.

Choose \(v_0,g_0\) with \(F(v_0)(g_0)\ne0\), and write the
adelic Iwasawa decomposition
\(g_0=u m_0 k_\infty x_f\).  Split
\(x_f=x^{(w)}\operatorname{loc}_w(x_w)\), with
\(x^{(w)}\) equal to one at \(w\).  The map
\[
 L(v)(m)=F(v)(m k_\infty x^{(w)})
 \tag{BCS.11}
\]
has values in the original tensor space of the twisted cuspidal
coefficients.  It is nonzero, since evaluation at
\(v=\sigma(x_w)v_0,m=m_0\) recovers the chosen nonzero value;
left unipotent invariance removes \(u\).  The element
\(k_\infty x^{(w)}\) commutes with the local group.  Writing a
local parabolic element as its Levi part times its unipotent part,
and moving the latter to the left, proves
\[
 L(\sigma(p)v)(m)=L(v)(m\,m_p),\qquad p\in P(F),
 \tag{BCS.12}
\]
where \(m_p\) is the original local Levi inclusion.

BCS4 makes the restriction of each twisted cuspidal space an
algebraic sum of copies of the local twist
\[
                         \rho\nu^{(a-1)s_{d,k}}.
\]
The multiplication map from the algebraic tensor of coefficient
spaces to their functions on the product of Levi groups is injective:
for each finite sum of pure tensors, independent evaluations on each
factor separate the finite-dimensional spans of its coordinates.
Thus the original \(\mathrm{tensorForms}\) is this actual tensor
space, not a further quotient.  It is therefore
isotypic for the tensor of these components.  Project to one copy
detecting a nonzero value of \(L\).  The resulting nonzero map
\(\ell\) satisfies (BCS.12) with values in this tensor.
Frobenius reciprocity is explicit here:
\[
                   \Phi(v)(g)=\ell(\sigma(g)v).
 \tag{BCS.13}
\]
It is smooth on the right, intertwines right translations, and has
the required unnormalized left covariance by (BCS.12).  The quotient
\(P(F)\backslash G(F)\) is compact, so there is no further
support condition.

For \(d\) equal blocks of size \(a\), the exponent of
\(|\det m_k|\) in \(\delta_P^{1/2}\) is
\[
 \frac{a(d-1-2k)}2=a s_{d,k}.
 \tag{BCS.14}
\]
Hence changing (BCS.13) to normalized induction subtracts
\(a s_{d,k}\) from its coefficient exponent
\((a-1)s_{d,k}\), leaving \(-s_{d,k}\).  This is exactly
(BCS.10).  Evaluation at the identity proves the map is nonzero.
All component maps, blocks and coefficient twists used here are the
ones from the original ascending embedding.
\end{proof}
% END BCS7

% BEGIN BCS8
\paragraph{BCS8. The precisely matched published generic-unitary input.}
\label{claim:BCS8}
For a nonarchimedean characteristic-zero local field, a smooth
irreducible generic unitarisable \(\rho\) has a presentation
\[
       \rho\simeq\delta_1\nu^{t_1}\times\cdots\times
                   \delta_b\nu^{t_b},\qquad |t_j|<\tfrac12,
 \tag{BCS.15}
\]
where the \(\delta_j\) are unitary square-integrable representations.
This is an explicit published input: M\oe glin--Waldspurger,
\emph{Le spectre r\'esiduel de \(\mathrm{GL}(n)\)},
\S\S I.10--I.11, pp.639--640,
\url{https://www.numdam.org/item/10.24033/asens.1595.pdf}.
Their \(\pi,\delta_j,v_j\) are respectively
\(\rho,\delta_j,t_j\); their normalized induction and absolute
value are (BCS.2).  The Whittaker hypothesis eliminates their longer
Speh factors.  All our \(E_w\) have characteristic zero.  BCS1 and
BCS5 supply precisely unitarisability and genericity for the actual
cuspidal local components, so (BCS.15) applies to each of them.
This application retains the strict interval; it does not assume
temperedness or prove a global residual-spectrum classification.
% END BCS8

% BEGIN BCS9
\paragraph{BCS9. The full unique cosocle of the descending Speh induction.}
\label{claim:BCS9}
For every \(\rho\) in BCS8 and every \(d\ge1\), the descending
product in BCS6 has an irreducible cosocle with a unique maximal
proper invariant subspace, in the literal sense of
\(\mathrm{IsCosocle}\).
\begin{proof}
Use the retained \texttt{SquareIntegrableSegment} to write
\(\delta_j=L(\Delta_j)\) for a segment.  The real center of each
\(\Delta_j\) is zero: its central character is the product of the
central characters of its segment terms, while unitarisability makes
that product have absolute value one.  More explicitly, if the first
term has rank \(r\), segment length \(l\), and central character
\(\omega\), then at a uniformizer \(z\)
\[
 \log|\omega(z)|=-\frac{r(l-1)}2\log|z|_F.
 \tag{BCS.16}
\]
Indeed each of the \(l\) terms contributes
\(\log|\omega(z)|+rh\log|z|_F\), for \(0\le h<l\),
and their sum is zero.  This also handles any unitary supercuspidal
line implicit in the segment notation.

If \(\Delta\nu^x\) precedes \(\Delta'\nu^y\) in the original
segment order and both underlying segments have unitary
square-integrable quotients, then
\[
                             x+1\le y.
 \tag{BCS.17}
\]
To check this exactly, express both on their common supercuspidal
line.  Precedence means their first endpoints differ by an integer
\(j\ge1\), their last endpoints differ by a positive integer,
and the intervals overlap or are adjacent.  If their lengths are
\(l,l'\), the last condition needed here is \(j+l'-l\ge1\).
Subtract (BCS.16) for the two first terms.  It gives
\(y-x=j+(l'-l)/2\), the average of the two positive endpoint
differences.  Therefore \(y-x\ge1\), proving (BCS.17).

Expand each factor of the descending product using (BCS.15), and
apply induction in stages.  The resulting ordered factors are
\[
          L\bigl(\Delta_j\nu^{t_j+s_{d,k}}\bigr),
       \qquad (k,j)=(0,1),\ldots,(0,b),(1,1),\ldots,(d-1,b).
 \tag{BCS.18}
\]
If an earlier segment preceded a later one, (BCS.17) would imply
\[
        t_j+s_{d,k}+1\le t_{j'}+s_{d,k'},\qquad k\le k'.
\]
But \(t_{j'}-t_j<1\) and
\(s_{d,k'}-s_{d,k}=k-k'\le0\), a contradiction.  Thus this is
exactly an ordered list of \(L\)-segment factors as required by
the original \texttt{StandardCosocleGL}.  That retained theorem
gives its irreducible unique cosocle.  Transporting through the
explicit induction-in-stages equivalence gives the asserted cosocle
of the original descending product.  Its uniqueness is the unique
maximal-kernel assertion of that theorem, so it applies to every
irreducible quotient, not merely one chosen Langlands quotient.
\end{proof}
% END BCS9

% BEGIN BCS10
\paragraph{BCS10. Every original discrete local component is the required Speh.}
\label{claim:BCS10}
Let \(V\) be discrete of positive rank and let
\(\operatorname{AscEmb}(V,P,d)\), with \(P\) normalized.
For every \(w\), every actual local component \(\rho\) of
\(P\), and every actual local component \(\sigma\) of \(V\),
\[
                         \mathrm{IsSpeh}(\rho,d,\sigma).
 \tag{BCS.19}
\]
\begin{proof}
BCS1 supplies the two positive invariant forms.  BCS7 supplies the
nonzero ascending map, and BCS6 gives a surjection from the
descending product onto \(\sigma\).  BCS5--BCS9 give the unique
cosocle of that product.  The kernel of a surjection onto an
irreducible representation is a maximal proper invariant subspace,
so it is the unique kernel of the cosocle map.  The induced quotient
equivalence identifies \(\sigma\) with this cosocle.  This proves
the literal \(\mathrm{IsSpeh}\), for every actual component map
in the statement and for every positive length.
\end{proof}
% END BCS10

% BEGIN BCS11
\paragraph{BCS11. The exact block Capelli shift.}
\label{claim:BCS11}
For the original unnormalized induced automorphic functions of block
ranks \(n_i\), with \(\sum_i n_i=N\), whose coefficient spaces
have infinitesimal characters \(m_i\), the infinitesimal character
on that induced space is
\[
       \sum_i\{z-\mu_i/2:z\in m_i\},\qquad
       \mu_i=\sum_{j>i}n_j-\sum_{j<i}n_j.
 \tag{BCS.20}
\]
\begin{proof}
Here is the algebraic block computation, including its sign.
For operators satisfying the \(\mathfrak{gl}_N\) relations, use
the source column determinant
\[
 q_N(X,u)=\sum_{\omega\in S_N}\operatorname{sgn}(\omega)
       \prod_{k=0}^{N-1}
       \bigl(X_{\omega(k),k}+\delta_{\omega(k),k}(u-k)\bigr),
 \tag{BCS.21}
\]
with factors in increasing column order.  Let \(J\) be the span
of \(X_{pq}A\), where the block of \(p\) is earlier than that of
\(q\).  It is a right ideal and is stable under left multiplication
by Levi entries.  The latter assertion follows immediately from
\([X_{ij},X_{pq}]=\delta_{jp}X_{iq}-\delta_{qi}X_{pj}\): the
commutator is again upper-block when \(i,j\) belong to one block.

For two consecutive blocks of sizes \(r,m\), the determinant
expansion gives
\[
 q_{r+m}(X,u)\equiv
       q_m(X_{22},u)q_r(X_{11},u-m)\pmod J.
 \tag{BCS.22}
\]
One way to see every discarded term is to move the last block before
the first simultaneously in the rows and columns.  The quantum
minor alternating identities preserve the determinant.  They follow
by interchanging adjacent columns in (BCS.21): the commutator of
their two entries is canceled by the difference of their consecutive
shifts; row interchange gives the ordinary alternating sign.  In the
new order, permutations preserving both blocks give exactly the
right side of (BCS.22).  Any other permutation has, in its initial
block of columns, an upper-block entry.  The entries preceding the
first such entry are Levi entries, and hence preserve \(J\) on the
left.  This term is in \(J\).  Entries in the two diagonal blocks
commute, so their two determinant factors may be interchanged.
Induction on blocks proves
\[
 q_N\!\left(X,u+\frac{N-1}{2}\right)
 \equiv\prod_i q_{n_i}\!\left(X_{ii},
           u+\frac{n_i-1}{2}-\frac{\mu_i}{2}\right)\pmod J.
 \tag{BCS.23}
\]
Indeed the shift for block \(i\) is also
\(u+(N-1)/2-\sum_{j>i}n_j\), which is the displayed expression.

Apply this to the right Lie operators on the source induced
functions.  At a Levi point, an upper-block operator kills a left
unipotent-invariant function: right infinitesimal multiplication
becomes left unipotent multiplication after conjugation by that
Levi point.  The same is true after any right differential operator,
since such operators preserve left unipotent invariance.  Thus every
element of \(J\) vanishes there.  On a pure tensor of coefficient
functions, each diagonal-block determinant acts by its specified
polynomial \(\prod_{z\in m_i}(u-\mu_i/2+z)\).  Linearity gives
the same conclusion on their algebraic tensor span.  Apply the
argument to every right translate by \(K_\infty G(\mathbb A_f)\);
adelic Iwasawa decomposition and left unipotent invariance then give
the identity at every point.  Right differentiation preserves these
induced coefficient spaces by the differentiated Iwasawa action,
so successive block determinants are legitimate.  This proves
(BCS.20) on the actual source functions.
\end{proof}
% END BCS11

% BEGIN BCS12
\paragraph{BCS12. The infinity character of the ascending embedding.}
\label{claim:BCS12}
Under the hypotheses of BCS10, for every embedding \(\tau\) there
is a multiset \(m_\tau\) of cardinality \(a\) such that
\[
 P.\mathrm{HasInfChar}(\tau,m_\tau),\qquad
 \mathrm{HasInfCharGL}\bigl(V,\tau,
                    \operatorname{spehInf}(d,m_\tau)\bigr).
 \tag{BCS.24}
\]
Any already supplied infinitesimal character of \(V\) equals this
one exactly.
\begin{proof}
Choose \(m_\tau\) by BCS4.  Twisting by
\(|\det|_{\mathbb A_E}^{t}\) adds \(t\) to each entry of
\(m_\tau\): differentiate the determinant character in each
diagonal Lie operator in (BCS.21).  For the equal blocks in (BCS.7),
\(\mu_k/2=a s_{d,k}\) by (BCS.14).  Therefore BCS11 gives on
the target of the ascending embedding the multiset
\[
             \sum_{k=0}^{d-1}\{z-s_{d,k}:z\in m_\tau\}.
\]
The involution \(k\mapsto d-1-k\) changes \(-s_{d,k}\) to
\(s_{d,k}\), so this is (BCS.3).  The ascending map intertwines
every Lie operator, hence every Capelli polynomial.  Its injectivity
transfers that scalar polynomial to \(V\).  The degree is
\(da=N\).  Uniqueness of the monic scalar polynomial in BCS4 proves
the final assertion; there is no unspecified common scalar twist.
\end{proof}
% END BCS12

% BEGIN BCS13
\paragraph{BCS13. The complete original \texttt{AscComp} consumer.}
\label{claim:BCS13}
The preceding assertions prove the literal \texttt{AscComp} of
\texttt{BaseChangeIdentities/Narrow.lean} for every totally complex
number field and every \(N\ge1\).
\begin{proof}
For the first clause, a normalized cusp of rank \(a=N\) has
\(V=P.V\) discrete by BCS1.  At length one the unique exponent
is zero, the parabolic is the whole group, and the identity inclusion
is the original \(\operatorname{AscEmb}(V,P,1)\).

For the second clause, take any discrete \(V\), normalized cusp
\(P\), and actual \(\operatorname{AscEmb}(V,P,d)\).
BCS4 provides both local components at every finite place.
BCS10 gives the \texttt{loc} clause for \emph{all} choices of
those components.  BCS12 gives the infinity clause.  These are
exactly the three fields \texttt{exLoc}, \texttt{loc}, \texttt{inf}
of \(\mathrm{HasSpehComp}(V,P,d)\).  This proves all of
\texttt{AscComp}; it does not assert the separate production of an
ascending embedding for an arbitrary discrete \(V\).
\end{proof}
% END BCS13

\subsection{The original spherical parameters and the final witnesses}

% BEGIN BCS14
\paragraph{BCS14. Exact fixed vectors of induction and quotients.}
\label{claim:BCS14}
For smooth local factors \(\tau_i\), evaluation at one identifies
\[
 (\tau_1\times\cdots\times\tau_c)^{K_n}
                  \simeq\bigotimes_i\tau_i^{K_{N_i}}.
 \tag{BCS.25}
\]
Compact fixed vectors are exact on smooth representations.  Hence a
spherical quotient of such an induction forces each nonzero-rank
factor to be spherical.
\begin{proof}
Iwasawa decomposition gives \(G=PK_n\).  A right \(K_n\)-fixed
section is determined by its value at one.  This value is fixed by
the integral Levi because \(\delta_P=1\) on \(P\cap K_n\).
Conversely, such a value defines
\(f(pk)=\delta_P(p)^{1/2}\tau(m_p)v\); two decompositions
differ by \(P\cap K_n\), so the definition is independent of
choices and is a smooth right-fixed section.

On the algebraic tensor space the average over the integral Levi is
the tensor product of the individual averages.  Their image is
exactly the tensor of the fixed spaces, proving (BCS.25).  An
algebraic tensor of nonzero vector spaces over a field is nonzero;
if some fixed space is zero, so is the tensor.

For a surjective smooth map \(T:X\to Y\) and \(y\in Y^K\),
lift \(y\) to \(x\) and average \(x\) over \(K\).  The finite
average has image \(y\).  Injection is automatic, proving exactness.
A nonzero fixed vector in a quotient therefore lifts to a nonzero
fixed vector in the induction, and (BCS.25) proves the assertion.
\end{proof}
% END BCS14

% BEGIN BCS15
\paragraph{BCS15. The normalized finite Satake calculation.}
\label{claim:BCS15}
For smooth irreducible spherical \(\tau_i\), choose their
original Satake multisets \(B_i\subset\mathbb C^\times\).
Then every nonzero fixed vector of their normalized product has
Satake multiset \(\sum_iB_i\).  The original normalization is
\[
 T_j\bigm|_{\tau^{K_r}}=
             q_w^{j(r-j)/2}e_j(B)\qquad(0\le j\le r),
 \tag{BCS.26}
\]
where \(T_j\) is the characteristic function of
\(K_r\operatorname{diag}(\varpi I_j,I_{r-j})K_r\), and
\(\operatorname{vol}(K_r)=1\).
\begin{proof}
We first record why these parameters exist for each irreducible
spherical factor.  The spherical Hecke algebra is commutative:
transpose is a Haar-preserving anti-involution fixing each double
coset, by elementary divisors.  It is generated by
\(T_1,\ldots,T_r,T_r^{-1}\).  To verify generation, write a
dominant integral elementary-divisor tuple as a sum of the tuples
\((1^j,0^{r-j})\) and a central tuple.  The corresponding product
of \(T_j\)'s has the indicated top double coset with nonzero
coefficient and only lower dominance tuples.  This follows by
successively taking the dimensions of the reductions of the
associated lattices.  There are finitely many lower tuples with
the fixed determinant and with minimum entry bounded below, so
induction eliminates them.  Inverses of \(T_r\) supply the central
negative tuples.

The fixed space of an irreducible spherical representation is a
simple module for this algebra.  If a fixed subspace generates the
representation under \(G\), averaging its translates back to
\(K\) gives its Hecke span; thus a nonzero Hecke submodule is the
whole fixed space.  Admissibility makes this space finite dimensional.
A simple finite-dimensional module over a commutative complex algebra
is one dimensional.  Its eigenvalues in (BCS.26) define the
elementary symmetric coefficients of a monic polynomial of degree
\(r\).  Its constant coefficient is nonzero because \(T_r\) is
invertible.  Its roots give a multiset in \(\mathbb C^\times\),
uniquely.  The same calculation works for every uniformizer, since
changing it by a unit does not change any of the double cosets.

Here is the block calculation that fixes the normalization.  For
\(\sum_i j_i=j\), decomposition of the finite flags for
\(K_n\operatorname{diag}(\varpi I_j,I_{n-j})K_n/K_n\) by their
intersections with the block flag gives the normalized constant term
\[
 T_j\longmapsto
 \sum_{\substack{0\le j_i\le N_i\\\sum_i j_i=j}}
 q_w^{\{j(n-j)-\sum_i j_i(N_i-j_i)\}/2}
                         \bigotimes_i T_{i,j_i}.
 \tag{BCS.27}
\]
For explicit counting, the off-diagonal graph coordinates contribute
\(q_w^{\sum_{i<l}j_i(N_l-j_l)}\).  The modulus correction in
normalized induction contributes
\(q_w^{-\frac12\sum_{i<l}(j_iN_l-j_lN_i)}\).
Adding these exponents gives exactly the exponent in (BCS.27).
The diagonal flags give the indicated \(T_{i,j_i}\); these cells
partition the finite quotient, so no multiplicity is omitted.

Under evaluation at one in BCS14, this constant term is the action
on the tensor fixed vector.  Substitute (BCS.26) for each factor.
All powers of \(q_w\) combine to \(q_w^{j(n-j)/2}\), and the
remaining sum is
\(\sum_{\sum j_i=j}\prod_i e_{j_i}(B_i)=e_j(\sum_iB_i)\).
The fixed space is one dimensional by (BCS.25); hence every nonzero
fixed vector has these eigenvalues.  This proves exactly the source
\(\mathrm{HasSatGL}\) condition for the full product.
\end{proof}
% END BCS15

% BEGIN BCS16
\paragraph{BCS16. The exact Speh strings in the inert parameter.}
\label{claim:BCS16}
If \(d\ge1\), \(\mathrm{IsSpeh}(\rho,d,\sigma)\), and
\(\mathrm{HasSatGL}(w,\sigma,B)\), then some multiset
\(S\subset\mathbb C^\times\) satisfies
\[
 \mathrm{HasSatGL}(w,\rho,S),\qquad
 B_{\mathbb C}=\sum_{k=0}^{d-1}
                       \{zq_w^{s_{d,k}}:z\in S\}.
 \tag{BCS.28}
\]
Consequently, if the full product of Speh factors has Satake multiset
\(T\), then
\[
             T_{\mathbb C}=\sum_i\sum_{k=0}^{d_i-1}
                       \{zq_w^{s_{d_i,k}}:z\in S_i\},
 \qquad\mathrm{HasSatGL}(w,\rho_i,S_i).
 \tag{BCS.29}
\]
Here subscript \(\mathbb C\) means the original map from units
to their values, retaining multiplicities.
\begin{proof}
Lift a nonzero fixed vector of \(\sigma\) through its actual
Speh quotient by BCS14.  Formula (BCS.25) makes each
\(\rho\nu^{s_{d,k}}\) spherical.  These unramified determinant
twists are trivial on \(K_a\), so \(\rho\) is spherical.
Choose its unique multiset \(S\) by BCS15.

Twisting by \(\nu^t\) multiplies the Satake roots by
\(q_w^{-t}\): on the \(j\)-th double coset the determinant has
absolute value \(q_w^{-j}\), so its Hecke operator is multiplied
by \(q_w^{-tj}\).  Apply BCS15 to the descending induction and
transport the eigenvalue equations through its quotient map.  Its
root multiset is therefore \(\sum_k q_w^{-s_{d,k}}S\).
The index involution \(k\mapsto d-1-k\) changes this to the right
side of (BCS.28).  The image vector is nonzero, so uniqueness of the
scalar polynomial proves equality with \(B\).  Apply BCS15 first
to the outer product and then (BCS.28) to each factor to get
(BCS.29).  In particular \(q_w\) is the absolute norm of \(w\),
not the separate parameter \(q\) appearing in
\(\mathrm{IsUnramInv}(\varpi,s,q)\).
\end{proof}
% END BCS16

% BEGIN BCS17
\paragraph{BCS17. Positive ranks and distinct equal-length cuspidal classes.}
\label{claim:BCS17}
For any \texttt{DefiniteN} witness and the ascending cusps it
supplies, write \(N_i=d_i a_i\).  Then \(a_i,d_i\ge1\),
\(c>0\), \(\sum_i d_i a_i=n\), and
\[
 d_i=d_j,\quad
 (a_i,[P_i])=(a_j,[P_j])\text{ in }\coprod_a\mathrm{CG}(E,a)
                         \quad\Longrightarrow\quad i=j.
 \tag{BCS.30}
\]
\begin{proof}
The rank identities follow directly from each ascending embedding;
positive \(N_i\) gives both positive factors.  If \(c=0\), the
rank sum would give \(n=0\), contradicting \(n\ge1\).

At any complex embedding \(\tau\), the original \(\pi\) has a
dominant integral weight \(\lambda\) satisfying
\[
                    \mathrm{HasWeightAt}(\pi,\tau,\lambda).
\]
To see existence
on these actual functions, take the nonzero finite-dimensional
infinity-orbit span of a vector in \(\pi\), restrict it to the
\(U(n)\) factor at \(\tau\), and average an inner product over
that compact group.  A simple constituent is one of the original
highest-weight modules by compact \(U(n)\) highest-weight theory.
Its inclusion is exactly the map in \(\mathrm{HasWeightAt}\),
because the factor inclusion is the original \(cEmb_\tau\)
on \(\operatorname{embU}(n,\tau)\).

The source multiset is
\[
 \operatorname{infOfWeight}(\lambda)=
       \{\lambda_k+(n-1)/2-k:0\le k<n\}.
 \tag{BCS.31}
\]
Its entries are pairwise distinct: for \(k<l\), their difference
is \(\lambda_k-\lambda_l+l-k>0\).  By BCS12 and the unchanged
base-change weight clause,
\[
 \sum_i\operatorname{spehInf}(d_i,m_{i,\tau})
                           =\operatorname{infOfWeight}(\lambda).
 \tag{BCS.32}
\]
Each summand has positive cardinality \(N_i\).  Equality of the
cuspidal classes in (BCS.30) identifies their infinitesimal
characters, by the all-operator equivalence and BCS4.  Equal lengths
then make the two summands in (BCS.32) equal.  If \(i\ne j\),
any entry of this nonempty multiset occurs at least twice in the
sum, contradicting (BCS.31).  This proves the original distinctness
clause without assuming a multiplicity-free base-change decomposition.
\end{proof}
% END BCS17

% BEGIN BCS18
\paragraph{BCS18. The entire original \texttt{BcSpeh}, with unchanged witnesses.}
\label{claim:BCS18}
Under the retained inputs stated at the start of this chapter, the
literal \texttt{BcSpeh} holds.  More precisely, for every original
\(E,n,\pi\) there are \(c>0\), \(a_i,d_i\ge1\),
\(\sum_i d_i a_i=n\), actual normalized unitary conjugate-self-dual
cusps \(P_i\), actual local cusp components \(\rho_{i,w}\),
multisets \(m_{i,\tau}\), and the original discrete local components
\(\sigma_{i,w}\), with all seven fields of \texttt{Parts}, both
local unitarisability assertions, the literal \(\mathrm{IsSpeh}\),
and the full split local-product component assertion.
\begin{proof}
Use BCS3 to obtain the original \texttt{DefiniteN} data.  From each
\(\mathrm{CsdAsc}(V_i)\), choose \(a_i,d_i,P_i\) and the
supplied ascending embedding.  The global normalization, unitarity
and conjugate self-duality of \(P_i\) are retained from that very
record.  BCS13 supplies its \(\mathrm{HasSpehComp}\).  Choose
\(\rho_{i,w}\) and \(m_{i,\tau}\) from its existence fields;
retain the \emph{original} \(\sigma_{i,w}\) from base change.
The universal local clause of BCS10 applies to those components.
BCS1 supplies the two local forms.  Replacing the equality
\(N_i=d_i a_i\) only relabels finite index sets by their canonical
rank equality; no new representation is substituted.

Here is every field of the target record.
\begin{center}
\begin{tabular}{p{0.16\linewidth}p{0.74\linewidth}}
\texttt{pos} & BCS17 gives \(a_i,d_i\ge1\) and the exact rank sum.\\
\texttt{csd} & The \(P_i\) retain their original conjugate self-duality.\\
\texttt{loc} & The chosen \(\rho_{i,w}\) have the actual cusp component maps.\\
\texttt{inf} & Each \(m_{i,\tau}\) is the actual cusp infinity character.\\
\texttt{weight} & BCS12 replaces every original \(M_{i,\tau}\) by
\(\operatorname{spehInf}(d_i,m_{i,\tau})\); the original weight
identity holds for every \(\tau,\lambda\) satisfying
\(\mathrm{HasWeightAt}\).\\
\texttt{distinct} & BCS17 proves precisely (BCS.30).\\
\texttt{inert} & Apply BCS16 to (BCS.6), obtaining (BCS.29) with
\(T=\operatorname{bcParam}(n,\alpha)\), for every original
\(w,s,\varpi,q,\alpha\) and every original guard.\\
\end{tabular}
\end{center}
The two unitarisability clauses and \(\mathrm{IsSpeh}\) were just
verified for every \(i,w\).  Finally, the split clause of BCS3 is
already \(\mathrm{IsCompAlong}\) for the full product of these
same \(\sigma_{i,w}\) along the same \(\operatorname{locU}\)
and \(cLoc_w\).  It is copied without alteration.  This is the
full existential statement \texttt{BcSpeh}, with all positive
lengths allowed and with the original rank \(n\).
\end{proof}
% END BCS18

% BEGIN BCS19
\paragraph{BCS19. All-unramified spherical cusp factors for the same witnesses.}
\label{claim:BCS19}
The witnesses in BCS18 also satisfy, simultaneously at every finite
place,
\[
 \mathrm{IsUnramU}(\pi,w)\quad\Longrightarrow\quad
           \rho_{i,w}^{K_{a_i}}\ne0\quad\text{for every }i.
 \tag{BCS.33}
\]
The assertion also holds for every alternative actual local component
of the same \(P_i\).
\begin{proof}
The stronger \texttt{DefiniteN} witness retained in BCS3 makes the
full product of the \(\sigma_{i,w}\) spherical at every such
place.  BCS14 makes each \(\sigma_{i,w}\) spherical.  Lift a
nonzero fixed vector through its actual descending Speh quotient
and apply BCS14 once more.  Each determinant twist of
\(\rho_{i,w}\) is then spherical, and the twists are trivial on
\(K_{a_i}\), proving (BCS.33).  BCS4 identifies any alternative
actual component with the chosen one, and an intertwining
equivalence transports its nonzero fixed vector.  This argument
requires neither a new choice of \texttt{BcSpeh} at different places
nor a split/inert case distinction.
\end{proof}
% END BCS19

\paragraph{Scope of the completed bridge.}
BCS3 proves the original \texttt{DefiniteN}; BCS13 proves the original
bounded \texttt{AscComp}; BCS18 proves the original \texttt{BcSpeh};
BCS19 supplies the additional all-unramified consequence for precisely
the same witnesses.  The named \texttt{BaseChangeIsobaricN} foundation
remains an input.  The separate global Galois-attachment record,
ordinary Hecke producers, arithmetic row counts, and terminal RH/GRH
obligations are not conclusions of this chapter.
